
de @@M0@@
aient des degrés compris entre @@M1@@ et @@M2@@. On procède alors par récurrence sur @@M3@@.

Lorsque @@M4@@, i.e. @@M5@@, soit @@M6@@ la valeur commune. L'homomorphisme canonique (1.2.1) :
@@M7@@
est alors un isomorphisme. En effet, la vérification est locale, et d'après (1.2.3), on peut alors supposer que @@M8@@, et le résultat est trivial.

Dans le cas @@M9@@, on considère encore l'homomorphisme (1.2.1) :
@@M10@@
On vérifie localement que cet homomorphisme est injectif, et que son conoyau @@M11@@ est localement libre de type fini, à générateurs locaux de degré compris entre @@M12@@ et @@M13@@ : en effet, au voisinage de tout point de @@M14@@, @@M15@@ est soit nul, soit la composante homogène non nulle de @@M16@@ dont le degré est minimum; utilisant une décomposition de @@M17@@ en facteurs directs du type (1.2.2), on voit dans ce dernier cas en utilisant (1.2.3) que @@M18@@ est isomorphe au facteur direct @@M19@@ de @@M20@@, d'où les propriétés annoncées. On applique alors l'hypothèse de récurrence à @@M21@@, d'où aussitôt la conclusion voulue.

1.8. Revenant à la démonstration du théorème 1.1, montrons que @@M22@@ est engendré comme @@M23@@-module par les puissances de @@M24@@.

Si @@M25@@ est un @@M26@@-Module localement libre de type fini, il existe d'après 1.4 un @@M27@@-Module gradué de présentation finie @@M28@@, plat sur @@M29@@, tel que @@M30@@. D'après 1.6, @@M31@@ admet une résolution graduée par des @@M32@@-Modules gradués localement libres de type fini. On en déduit une
suite exacte :
